\RequirePackage{pdfmanagement}
\SetKeys[document/metadata]{pdfversion=1.7, pdfstandard=A-2u, lang=en-GB}
\input{preamble/pdfA-compliance/glyphtounicode}
\documentclass[11pt,a4paper]{report}
\newcommand*{\PaperTitle}{The complexity of multistage minimum utility cost}
\newcommand*{\PaperTitlePlain}{The complexity of multistage minimum utility cost}
\newcommand*{\PaperShortTitle}{Multistage minimum utility cost}
\newcommand*{\PaperSubtitle}{chemistry-open-problems \textperiodcentered{} 5 November 2026 \textperiodcentered{} preliminary}
\newcommand*{\PaperCredits}{Solver and manuscript: GPT-6-Astra (effort xhigh) \textperiodcentered{} Deliverables completed: Claude Opus 5.5 (effort xhigh) \textperiodcentered{} Independent verification: pending}
\newcommand*{\PaperKeywords}{heat exchanger networks, utility minimization, strong NP-hardness, 3-Partition, stream mixing}
\input{chemop-preamble}
\addbibresource{references.bib}
\newcommand{\CC}{\mathsf{C}}
\newcommand{\BB}{\mathsf{B}}
\newcommand{\NN}{\mathsf{N}}
\newcommand{\ER}{\exists\mathbb{R}}
\newcommand{\pos}[1]{\left(#1\right)_+}
\begin{document}
\makepapertitle
\begin{paperabstract}
Letsios et al. ask whether minimum utility cost in a multistage heat-exchanger
superstructure is polynomial-time solvable or NP-hard. We distinguish their
displayed journal equations, which omit temperature-approach constraints, from
two physical completions: approaches at the mixed stage boundaries and at the
actual counter-current exchanger endpoints. The displayed equations have a
closed-form optimum. Both physical completions are strongly NP-hard, even with
exactly two stages, all stream pairs permitted, unit utility prices and the
question of zero utility. A quadratic temperature identity forces every mixing
variance and every exchanger approach slack to vanish in the reduction. This
makes each hot stream choose a bin in a 3-Partition instance. Padding proves
hardness for every fixed number of stages at least two. With one stage the
counter-current version is an LP, whereas the boundary version is strongly
NP-complete. The boundary decision problem is in NP; the counter-current
problem is in the existential theory of the reals, and we do not establish its
general membership in NP. Exact symbolic and rational checks corroborate the
proofs; floating-point optimization is not used as a certificate.
\end{paperabstract}

\section*{The question and the model correction}
Dimitrios Letsios, Radu Baltean-Lugojan, Francesco Ceccon, Miten Mistry,
Johannes Wiebe and Ruth Misener, \emph{Approximation algorithms for process
systems engineering}, \emph{Computers \& Chemical Engineering} 132 (2020),
106599, DOI 10.1016/j.compchemeng.2019.106599, state in Section 6.4,
journal page 10 \autocite{Letsios2020}:
\begin{openquote}
``The complexity of the multistage minimum utility cost problem is an intriguing open question.''
\end{openquote}
The supplied journal PDF was checked visually. Its inequalities (5j,k) have the
correct decreasing signs, unlike arXiv v1. Equation (5m) still indexes cold
streams by the hot set; we necessarily read that index as the cold set.
No constraint in the display compares a hot temperature to a cold temperature.
Thus a physical interpretation requires an explicit additional assumption.
Our answer is for the two completions defined below, not a claim that the
missing constraints were printed in the source.

\begin{theorem}\label{thm:main}
For each fixed integer $r\geq2$, deciding whether the minimum utility cost is
zero in the $r$-stage counter-current model $\CC$ is strongly NP-hard.
This holds with positive integer stream capacities, all hot--cold pairs
available, both utility prices equal to one, and a minimum approach fixed to
either zero or any prescribed nonnegative rational constant. Temperatures
can be positive integers when the fixed approach is an integer, and have
constant denominators for any other fixed rational approach. The same
restrictions give strong NP-hardness
for the stage-boundary model $\BB$. Consequently, both problems are strongly
NP-hard when the stage count is part of the input, in unary or binary.

For $\BB$, the decision problem with any rational cost threshold is in NP
and hence strongly NP-complete. For $\CC$ it is in $\ER$; no general NP
membership assertion is made. For one stage, $\CC$ is polynomial-time
solvable and $\BB$ is strongly NP-complete.
\end{theorem}

\section{Exact formulation and input encoding}
Let $I,J$ be nonempty hot and cold stream sets, of sizes $n,s$. The given
positive heat-capacity flow rates are $F_i,G_j$. Given temperatures satisfy
$H_i\geq\underline H_i\geq0$ and $\overline C_j\geq C_j\geq0$:
hot stream $i$ must cool from $H_i$ to $\underline H_i$, and cold stream $j$
must heat from $C_j$ to $\overline C_j$. The stage count is $\ell\geq1$.
Unless explicitly fixed, it is an input integer. All numerical data, including
the common utility prices $c_{\rm CU},c_{\rm HU}\geq0$, minimum approach
$\delta\geq0$, and decision threshold, are binary-encoded rationals.
Our main reductions fix $\ell$ and both prices; they do not encode complexity
in a large rational denominator or in the number of stages.

Index boundaries by $k=0,\ldots,\ell$. Hot streams move from boundary 0 to
$\ell$, and cold streams move in the opposite direction. Their mixed
temperatures are $h_{ik},c_{jk}$. Stage $k\in\{1,\ldots,\ell\}$ has variables
$f_{ijk},g_{ijk}$ for the hot and cold branch capacities, $a_{ijk},b_{ijk}$
for their respective outlet temperatures, and $q_{ijk}$ for exchanged heat.
These are real decision variables, all nonnegative. Branches split from one
mixed inlet and mix again before the next stage; there is no cross-stream
material mixing. All pairs $(i,j)$ are available in every stage.
The base equations, denoted $\NN$, are the journal's (5), with only the
necessary cold-index normalization and notation changes:
\begin{align}
 h_{i0}&=H_i,& h_{i\ell}&\geq\underline H_i,&
 c_{j\ell}&=C_j,&c_{j0}&\leq\overline C_j,\label{eq:ends}\\
 h_{i,k-1}&\geq h_{ik},&&c_{j,k-1}\geq c_{jk},\label{eq:mono}\\
 \sum_j f_{ijk}&=F_i,&&\sum_i g_{ijk}=G_j,\label{eq:split}\\
 q_{ijk}&=f_{ijk}(h_{i,k-1}-a_{ijk})
          =g_{ijk}(b_{ijk}-c_{jk}),\label{eq:heat}\\
 F_i h_{ik}&=\sum_j f_{ijk}a_{ijk},&&
 G_j c_{j,k-1}=\sum_i g_{ijk}b_{ijk}.\label{eq:mix}
\end{align}
The nonnegative end utilities and objective are
\begin{equation}\label{eq:obj}
 U_i=F_i(h_{i\ell}-\underline H_i),\qquad
 V_j=G_j(\overline C_j-c_{j0}),\qquad
 \min\ c_{\rm CU}\sum_i U_i+c_{\rm HU}\sum_j V_j.
\end{equation}
Zero-heat branches may carry bypass flow. There is no minimum positive heat
duty, match-count restriction, exchanger-area bound, pressure-drop term,
capital charge, or intermediate utility. End utilities are available without
capacity restrictions. These assumptions isolate the source's utility-cost
problem from total annual cost and minimum matches.

Stages here are ordered structural positions, not prescribed temperature
intervals. Their boundary temperatures are optimized. In particular, the
fixed temperature-interval transshipment model discussed separately in the
source is not being substituted for this model.

\subsection{The two physical completions}
Model $\CC$ adds, for every $q_{ijk}>0$, the counter-current endpoint constraints
\begin{equation}\label{eq:C}
 h_{i,k-1}\geq b_{ijk}+\delta,\qquad
 a_{ijk}\geq c_{jk}+\delta.
\end{equation}
Model $\BB$ instead adds, for every $q_{ijk}>0$, the stage-boundary constraints
\begin{equation}\label{eq:B}
 h_{i,k-1}\geq c_{j,k-1}+\delta,\qquad
 h_{ik}\geq c_{jk}+\delta.
\end{equation}
An implication is imposed only on heat-carrying matches; it does not forbid
unheated bypass branches between temperature-incompatible streams.
Yee--Grossmann SYNHEAT uses the stage-boundary version. The public GAMS
implementation explicitly activates both boundary comparisons with its match
binary and sets the approach variable's lower bound to the prescribed minimum
\autocite{Yee1990,GAMS}. For non-isothermal mixing the actual exchanger outlets
are the relevant endpoints. The contemporary OpenHENS implementation displays
exactly those two branch comparisons, alongside its isothermal boundary
alternative \autocite{OpenHENS}. We use these sources to fix the approach
constraints, not their additional area objectives or design restrictions.

Adding $\delta$ to every cold temperature transforms either set of approach
constraints to its zero-approach version and leaves the heat balances
unchanged. This is the usual temperature shift mentioned in the source.
Conversely, all our hardness instances for $\delta=0$ become instances for any
fixed $\delta>0$ by adding $\delta$ to every hot temperature. A further common
translation makes all absolute temperatures strictly positive.

Summing the branch equations gives the stage balances
\begin{equation}\label{eq:stage}
 R_{ik}:=\sum_jq_{ijk}=F_i(h_{i,k-1}-h_{ik}),\qquad
 S_{jk}:=\sum_iq_{ijk}=G_j(c_{j,k-1}-c_{jk}).
\end{equation}
Write $E_H=\sum_iF_i(H_i-\underline H_i)$,
$E_C=\sum_jG_j(\overline C_j-C_j)$ and $Q=\sum_{ijk}q_{ijk}$.
The objective equals
\begin{equation}\label{eq:Qobj}
 c_{\rm CU}E_H+c_{\rm HU}E_C-(c_{\rm CU}+c_{\rm HU})Q.
\end{equation}

\begin{lemma}[note=Isothermal lifting]\label{lem:lift}
Any nonnegative $q,h,c$ satisfying the endpoint and stage balances can be
lifted to $\NN$. If they satisfy the active boundary comparisons, they can be
lifted simultaneously to $\BB$ and $\CC$.
\end{lemma}
\begin{proof}
In each stage set $a_{ijk}=h_{ik}$ and $b_{ijk}=c_{j,k-1}$.
When $R_{ik}>0$ set $f_{ijk}=F_iq_{ijk}/R_{ik}$; when $R_{ik}=0$ use any
nonnegative row summing to $F_i$. Likewise, when $S_{jk}>0$ set
$g_{ijk}=G_jq_{ijk}/S_{jk}$, and otherwise use any column summing to $G_j$.
Equations \eqref{eq:heat} and \eqref{eq:mix} follow immediately from
\eqref{eq:stage}. On active matches, \eqref{eq:C} becomes \eqref{eq:B}.
\end{proof}
Thus $\BB$ and $\CC$ are compared through their feasible $(q,h,c)$ projections;
an arbitrary non-isothermal assignment satisfying $\BB$ need not itself
satisfy $\CC$. This distinction is important in applying the next proof to
$\BB$.

\section{The displayed journal equations have a closed-form optimum}
\begin{theorem}\label{thm:literal}
For every $\ell\geq1$, the optimum of $\NN$ is
\begin{equation}\label{eq:literal}
 c_{\rm CU}\max(E_H-E_C,0)+c_{\rm HU}\max(E_C-E_H,0).
\end{equation}
An optimal rational solution can be constructed in polynomial time, using
only one nonidle stage.
\end{theorem}
\begin{proof}
The stage balances imply $Q\leq\min(E_H,E_C)$, giving the lower bound in
\eqref{eq:literal}. Greedily transport heat from hot supplies
$F_i(H_i-\underline H_i)$ to cold demands $G_j(\overline C_j-C_j)$, with no
pair restrictions, until one side is exhausted. Let $q_{ij}$ be this
transportation matrix, $R_i$ its row sums and $S_j$ its column sums. In stage 1
set $h_{i1}=H_i-R_i/F_i$, $c_{j0}=C_j+S_j/G_j$ and $c_{j1}=C_j$.
Use the lift in \zcref{lem:lift}; repeat these post-stage temperatures through
all later, idle stages. Then $Q=\min(E_H,E_C)$. All entries have polynomial
rational bit length. A binary-encoded large $\ell$ is represented by this
one active stage and repeated idle stages, rather than exponentially many
explicit entries.
\end{proof}
For example, one capacity-one hot stream $1\to0$ and one capacity-one cold
stream $10\to11$ have zero utility in $\NN$: the exchanger transfers one heat
unit uphill in temperature. Thus \eqref{eq:literal} is a model correction,
not a resolution of the intended physical question.

\section{A quadratic identity that removes fractional mixing}
We first take $\delta=0$. For a single exchanger abbreviate its hot inlet and
outlet to $h,a$ and its cold outlet and inlet to $b,c$. Its balances imply
\begin{equation}\label{eq:single}
 f(h^2-a^2)-g(b^2-c^2)=q\bigl[(h-b)+(a-c)\bigr].
\end{equation}
For each stage define the nonnegative mixing variances
\begin{equation}\label{eq:var}
 W^H_k=\sum_{ij}f_{ijk}(a_{ijk}-h_{ik})^2,\qquad
 W^C_k=\sum_{ij}g_{ijk}(b_{ijk}-c_{j,k-1})^2.
\end{equation}
\begin{lemma}[note=Quadratic budget and equality case]\label{lem:quad}
Every feasible $\CC$ solution satisfies
\begin{align}
 &\sum_i F_i(H_i^2-h_{i\ell}^2)
 -\sum_j G_j(c_{j0}^2-C_j^2)\notag\\
 &\quad=\sum_{kij}q_{ijk}
 \bigl[(h_{i,k-1}-b_{ijk})+(a_{ijk}-c_{jk})\bigr]
 +\sum_k(W^H_k+W^C_k)\ \geq0.\label{eq:budget}
\end{align}
If the left side is zero, every positive-flow branch exits at its stream's
mixed outlet temperature, and every active exchanger satisfies
\begin{equation}\label{eq:interval}
 h_{i,k-1}=c_{j,k-1},\qquad h_{ik}=c_{jk}.
\end{equation}
\end{lemma}
\begin{proof}
Weighted variance and \eqref{eq:mix} give
$\sum_j f_{ijk}a_{ijk}^2=F_i h_{ik}^2+
\sum_j f_{ijk}(a_{ijk}-h_{ik})^2$, with the analogous cold identity.
Substitute these in the sum of \eqref{eq:single}. The difference of mixed
hot and cold quadratic drops in each stage equals the exchanger terms plus
both variances. Summing stages telescopes to \eqref{eq:budget}.
Every term on the right is nonnegative by \eqref{eq:C}; for $q=0$ the
exchanger term is zero regardless of unused temperatures. Zero total forces
zero variances. For $q>0$, both branch capacities are positive by
\eqref{eq:heat}, so $a=h_{ik}$ and $b=c_{j,k-1}$; both nonnegative approach
slacks also vanish. These statements imply \eqref{eq:interval}.
\end{proof}
The identity is algebraic; it is not an unproved appeal to entropy or an
assumption of isothermal mixing. It forces that mixing property on the
particular instances constructed next.

\section{Strong NP-hardness with two stages}
Start from 3-Partition: $3m$ positive integers $a_1,\ldots,a_{3m}$ and $B$
satisfy
\begin{equation}\label{eq:3p}
 B/4<a_i<B/2,\qquad \sum_i a_i=mB.
\end{equation}
The question is whether they can be partitioned into $m$ triples of sum $B$.
This is strongly NP-complete \autocite{GareyJohnson}.
Set $T=m+1$, $\ell=2$, $\delta=0$, both utility prices to one, and the cost
threshold to zero. Introduce the following streams:
\begin{center}
\begin{tabular}{llll}
\hline
Stream & Capacity & Inlet & Target\\
\hline
Hot item $i$ & $a_i$ & $T$ & $0$\\
Cold upper $U_j$, $1\leq j\leq m$ & $B$ & $j$ & $T$\\
Cold lower $V_j$, $1\leq j\leq m$ & $B$ & $0$ & $j$\\
\hline
\end{tabular}
\end{center}
There are no forbidden matches or bounds on their number. Energy and
quadratic temperature budgets agree exactly:
\begin{align}
 E_H=E_C&=mBT,\label{eq:equalenergy}\\
 \sum_i a_iT^2&=\sum_{j=1}^m B\bigl[(T^2-j^2)+j^2\bigr]=mBT^2.
 \label{eq:equalsquare}
\end{align}

\begin{lemma}\label{lem:reduction}
The constructed $\CC$ instance has zero utility if and only if the
3-Partition instance is a YES instance. The equivalence also holds for $\BB$.
\end{lemma}
\begin{proof}
Suppose first that $\CC$ has zero utility. Positive utility prices force every
stream to reach its prescribed target, so \eqref{eq:equalsquare} makes the
left side of \eqref{eq:budget} zero. Every active match consequently obeys
\eqref{eq:interval}.

No $U_j$ can receive heat in stage 2: its inlet there is $j>0$, whereas every
hot stream's outlet there is 0. These cannot be equal as required by
\eqref{eq:interval}. Hence $c_{U_j,1}=j$ and the entire heat $B(T-j)$ is
received in stage 1. Similarly, no $V_j$ can receive heat in stage 1:
its outlet there is its target $j<T$, whereas every hot inlet there is $T$.
Thus $c_{V_j,1}=j$ and its entire heat $Bj$ is received in stage 2.

Let $x_i=h_{i1}$. If $x_i<T$, hot stream $i$ releases positive heat in stage 1
and must meet some $U_j$, forcing $x_i=j$. If instead $x_i=T$, it releases
positive heat in stage 2 and must meet some $V_j$, forcing $T=j$, impossible.
Therefore every $x_i$ belongs to $\{1,\ldots,m\}$. For $x_i=j$, the only
possible active partners are $U_j$ in stage 1 and $V_j$ in stage 2.
The stage-1 balance of $U_j$ now gives
\begin{equation}\label{eq:bins}
 (T-j)\sum_{i:x_i=j}a_i=B(T-j),\qquad
 \sum_{i:x_i=j}a_i=B.
\end{equation}
The strict bounds in \eqref{eq:3p} force exactly three items in every group.
No fractional assignment to distinct bins survives the equality case.

Conversely, let $S_1,\ldots,S_m$ be a partition with capacity $B$ per group.
For $i\in S_j$ set its hot temperatures to $(T,j,0)$, exchange
$a_i(T-j)$ with $U_j$ in stage 1 and $a_i j$ with $V_j$ in stage 2.
Cold temperatures are $(T,j,j)$ for $U_j$ and $(j,j,0)$ for $V_j$.
These are exact stage balances and satisfy both boundary comparisons with
equality. The isothermal lift supplies every branch flow and outlet
temperature, including idle cold bypasses. Thus both $\BB$ and $\CC$ have
zero utility. Finally, any zero-utility $\BB$ solution can be lifted to a
$\CC$ solution with the same stage heat and temperatures by
\zcref{lem:lift}, proving the reverse implication for $\BB$ as well.
\end{proof}
The transformation uses $3m$ hot streams and $2m$ cold streams; its largest
temperature is $m+1$ and its capacities are $a_i$ or $B$. It is polynomial
and preserves polynomial bounds on numerical magnitudes, proving strong
NP-hardness. On this restricted family the zero-utility decision problem is
itself strongly NP-complete, since a bin partition is a polynomial witness.
This last statement does not establish NP membership for arbitrary $\CC$
instances.

\subsection{Every fixed larger number of stages}
Fix $r\geq2$, put $p=r-2$ and $A=mB$. Raise every hot inlet to $T+p$ and add
the following $p$ cold streams, leaving $U_j,V_j$ unchanged:
\begin{equation}\label{eq:padding}
 D_t:\quad \hbox{capacity }A,\quad T+t-1\ \longrightarrow\ T+t,
 \qquad t=1,\ldots,p.
\end{equation}
Use exactly $r$ stages. Their energy and squared-temperature increments
telescope, so both equalities in \eqref{eq:equalenergy}--\eqref{eq:equalsquare}
remain true with the new hot inlet $T+p$.
In a zero-utility solution, every positive hot temperature drop must exactly
coincide with an interval of some cold stream by \eqref{eq:interval}.
It therefore lies within that cold stream's full inlet-to-target interval.
No cold interval crosses any of the boundaries
$T,T+1,\ldots,T+p-1$. Each hot stream needs at least one nonzero stage in
each of the $p$ upper bands to descend to $T$. It needs at least two more
stages to descend from $T$ to 0, because no cold stream spans that whole
interval. There are only $p+2$ stages. Hence every hot stream uses exactly
one full drop in each upper band, in decreasing order, followed by exactly
the two stages of the original construction. There are no idle or additional
partial drops. The argument of \zcref{lem:reduction} applies to these last two stages.
Conversely all hot streams can traverse each added band together, supplying
the capacity-$A$ dummy cold stream, before executing the partition witness.
This proves hardness for each fixed $r$; adding $\delta$ to hot temperatures
proves it for any fixed nonnegative approach. With $\delta>0$ the witness
has genuinely positive physical endpoint differences, so the reduction does
not require a zero-driving-force exchanger in physical temperature coordinates.

\section{One stage distinguishes the two models}
\subsection{A polynomial LP for the counter-current version}
For $\ell=1$, let $D_{ij}=H_i-C_j-\delta$. Delete pairs with $D_{ij}\leq0$.
The exact feasible heat projection of $\CC$ is
\begin{align}
 q_{ij}&\geq0,&
 \sum_jq_{ij}&\leq F_i(H_i-\underline H_i),&
 \sum_iq_{ij}&\leq G_j(\overline C_j-C_j),\label{eq:LPheat}\\
 \sum_{j:D_{ij}>0}\frac{q_{ij}}{D_{ij}}&\leq F_i,&&
 \sum_{i:D_{ij}>0}\frac{q_{ij}}{D_{ij}}\leq G_j.\label{eq:LPflow}
\end{align}
Indeed, \eqref{eq:C} and \eqref{eq:heat} imply $q_{ij}\leq f_{ij}D_{ij}$
and $q_{ij}\leq g_{ij}D_{ij}$. Summing proves necessity. Conversely start
both branch capacities at $q_{ij}/D_{ij}$ on retained pairs and at zero
elsewhere. Allocate unused hot capacity arbitrarily within each row and
unused cold capacity within each column, independently, to meet their split
sums. On positive flow define $a_{ij}=H_i-q_{ij}/f_{ij}$ and
$b_{ij}=C_j+q_{ij}/g_{ij}$; for zero flow choose its inlet temperature.
The capacity lower bounds guarantee \eqref{eq:C}, and the mixed temperatures
and utility bounds follow from \eqref{eq:LPheat}. Maximize $\sum q_{ij}$ in
this rational LP, or use \eqref{eq:Qobj} directly. This proves polynomial
solvability, including construction of an optimal network.

\subsection{Strong NP-hardness of one-stage boundary approaches}
Use the same 3-Partition data. Give each hot stream capacity $a_i$ and range
$m\to0$; give cold stream $j=0,\ldots,m-1$ capacity $B$ and range $j\to m$.
Then $E_H=m^2B$, $E_C=Bm(m+1)/2$. At unit prices the threshold
$L_0=E_H-E_C=Bm(m-1)/2$ is attained only if all cold demands are met.

Suppose $\BB$ attains this threshold. For each hot stream with positive heat,
let $r_i$ be the largest cold inlet index among its active partners; put
$r_i=m$ if it is idle. Let $x_r=\sum_{i:r_i=r}a_i$. Boundary feasibility
forces its hot outlet to be at least $r_i$, so total recovered heat is at
most $\sum_{r<m}(m-r)x_r$. Consequently
\begin{equation}\label{eq:loss}
 L:=\sum_{r=0}^m r x_r\leq E_H-E_C=L_0.
\end{equation}
For each $k=1,\ldots,m-1$, all cold streams with inlet at least $k$ must draw
from hot streams classified $r\geq k$. Their demand and an upper bound on
their supply are
\begin{equation}\label{eq:suffix}
 D_k=B\sum_{j=k}^{m-1}(m-j),\qquad
 S_k=\sum_{r=k}^{m-1}(m-r)x_r,\qquad S_k\geq D_k.
\end{equation}
Use strictly positive multipliers
\begin{equation}
 \lambda_k=\frac{k}{m-k}-\frac{k-1}{m-k+1}
           =\frac{m}{(m-k)(m-k+1)}.
\end{equation}
Telescoping gives
$\sum_k\lambda_k S_k=\sum_{r=1}^{m-1}r x_r$ and
$\sum_k\lambda_k D_k=L_0$. Together with \eqref{eq:loss}, these imply
$x_m=0$ and equality in every suffix constraint. Successive differences of
\eqref{eq:suffix} give $x_r=B$ for $r=1,\ldots,m-1$; total capacity gives
$x_0=B$. Thus the items are partitioned into $m$ bins of capacity $B$ and
hence into triples. Conversely, assigning each partition group entirely to
its cold stream attains the threshold and obeys the boundary approaches.
The construction is strong and uses one stage.

These models really differ. For example hot $(F,H,\underline H)=(1,10,0)$
and cold streams $(G,C,\overline C)=(1,8,10),(1,0,8)$ have one-stage
maximum recovery $42/5$ in $\CC$, but only 8 in $\BB$, at $\delta=0$.
For $\CC$, the LP inequalities $q_1/2+q_2/10\leq1$ and $q_2\leq8$ give
the sharp bound $q_1+q_2\leq42/5$, attained at $(2/5,8)$.
For $\BB$, positive heat to the high cold stream forces the mixed hot
outlet to be at least 8, so total recovery is at most 2; omitting that match
permits exactly 8. With two stages both attain 10, by traversing $10\to8\to0$.

\section{Decision complexity and the encoding of the stage count}
For fixed or unary $\ell$, the projected $\BB$ problem is a union of
polynomial-size rational polyhedra: guess the set of permitted active matches,
fix $q=0$ off that set, and impose \eqref{eq:B} on it. Zero heat on a guessed
match is harmless. The remaining equations are \eqref{eq:stage}, bounds and
the cost threshold. A nonempty rational polyhedron has a rational feasible
point with polynomial bit length, giving NP membership. The lifting lemma
recovers the original branch variables without increasing that bound
superpolynomially. The reductions therefore give strong NP-completeness.

For $\CC$, the active implications can instead be written as polynomial
inequalities
\begin{equation}
 q_{ijk}(h_{i,k-1}-b_{ijk}-\delta)\geq0,\qquad
 q_{ijk}(a_{ijk}-c_{jk}-\delta)\geq0.
\end{equation}
Together with the bilinear equations, these give an existential formula over
the reals of polynomial size. This proves membership in $\ER$ for fixed or
unary stage counts. Bilinearity alone is not a proof of NP membership, and
we do not claim that all feasible instances have small rational witnesses.

For completeness the following stage bound removes an exponential expansion
when $\ell$ is encoded in binary. It uses the absence of forbidden matches
and of extra design restrictions in our precise model.
\begin{lemma}[note=Compression of the number of stages]\label{lem:compress}
Any terminal mixed temperatures achievable by $\CC$ or $\BB$ with any number
of stages are achievable by both models with at most $2(n+s)-1$ stages.
\end{lemma}
\begin{proof}
Shift cold temperatures by $\delta$ to take $\delta=0$; denote the attained
hot ends by $x_i$ and cold ends by $y_j$. Put $\phi_\theta(t)=\max(t-\theta,0)$.
Every feasible solution necessarily satisfies, for all real $\theta$,
\begin{equation}\label{eq:cut}
 A(\theta):=\sum_iF_i[\phi_\theta(H_i)-\phi_\theta(x_i)]
 \ \geq\
 D(\theta):=\sum_jG_j[\phi_\theta(y_j)-\phi_\theta(C_j)].
\end{equation}
To see this for $\CC$, an active exchanger has $a\geq c$ and $h\geq b$.
The uniform points $(1-u)a+uh$ dominate $(1-u)c+ub$ for $0\leq u\leq1$.
Since $\phi_\theta'$ is nondecreasing, the average derivative on $[a,h]$
is at least that on $[c,b]$. Multiplying by $q=f(h-a)=g(b-c)$ gives
$f[\phi_\theta(h)-\phi_\theta(a)]\geq
g[\phi_\theta(b)-\phi_\theta(c)]$.
The zero-heat case contributes zero. Convexity and the mixing equations add
nonnegative Jensen gaps for both hot and cold mixing, exactly as in
\zcref{lem:quad}. Sum and telescope to obtain \eqref{eq:cut}.
For $\BB$ use the isothermal lift first. Total recovered hot and cold heats
are equal by energy conservation.

Sort the distinct numbers among $H_i,x_i,y_j,C_j$ in decreasing order.
There are at most $2(n+s)$ such numbers. Use one stage per consecutive pair
$u>v$. Prescribe each cold stream's boundaries in that stage by clipping
$u,v$ to its range $[C_j,y_j]$. Thus any cold stream active in this interval
has outlet $u$, inlet $v$, and demand $G_j(u-v)$.

Process the intervals downwards. Before the interval $(v,u)$, a hot stream
with $H_i\geq u$ has current temperature $p_i\geq\max(x_i,u)$, by induction;
all other hot streams remain at their inlets. The available capacity to
release heat in this interval is
\begin{equation}
 F_i[p_i-\max(x_i,v)]
\end{equation}
for the former streams, and zero for the latter. The total of these
capacities is $A(v)-D(u)$: earlier intervals have consumed exactly $D(u)$,
and streams not yet eligible have no supply above $v$ because no inlet
lies strictly between consecutive endpoints. By \eqref{eq:cut}, it is at
least $D(v)-D(u)$, the current cold demand. Choose any transportation of
that demand within the available row capacities. Lower each hot temperature
by its allocated heat divided by $F_i$. On each active match its previous
temperature is at least $u$ and its next temperature is at least $v$, so
both boundary approaches hold. This also proves the induction invariant
for the next interval.

At the bottom all cold streams have reached their required $y_j$. Equality
of total hot and cold heat, and the lower bounds $p_i\geq x_i$, force every
hot stream to reach $x_i$. Apply the isothermal lift in each stage. The
construction uses at most $2(n+s)-1$ stages and proves the claim. If no
temperature interval exists, one idle stage suffices.
\end{proof}
Idle stages can always be inserted. Therefore, for the decision problem we
may replace a binary-encoded $\ell$ by
$\min\{\ell,2(n+s)-1\}$. If $\ell$ is below this bound nothing changes;
otherwise \zcref{lem:compress} preserves exactly the attainable terminal
temperatures and costs. The NP and $\ER$ memberships above consequently
hold for binary stage encoding too. This compression does not make two
stages equivalent to an unrestricted number of stages: the explicit example
above and the reductions retain their stage restrictions.

For these models minima are attained. Bound all boundary temperatures by
their stream endpoints. All branch flows and heats are then bounded.
In $\CC$, an active branch temperature lies between nonnegative stream
inlets and a global maximum inlet/target temperature by the approaches;
when $q=0$, any positive flow has its inlet temperature, and a zero-flow
temperature can be set to that inlet. Hence one may restrict all branch
temperatures to a compact box. The polynomial implication formulation is
closed. For $\BB$ the bounded projected feasible set is a finite union of
closed polyhedra and the lifting lemma applies. In particular, the
zero-utility decision in the reduction is about an attained optimum, not
an unattained infimum.

\section{Exact checks, exploration and limitations}
The proofs above stand independently of all computation.
\texttt{scripts/verify\_reduction.py} uses SymPy polynomial expansion and
Python rational arithmetic to check the exchanger and mixing identities,
the endpoint budgets, 36 full witnesses with stages 2--5 and approaches
$0,1,3/2$, five explicit equality-case enumerations, and every sorted two-bin
3-Partition promise instance with $7\leq B\leq16$ (21 instances, 18 YES and
3 NO). All equations, nonnegativity and both active approach versions are
checked on each saved witness. The negative finite enumerations use the
equality lemma to reduce the feasible paths to a finite set; they are not
independent nonlinear infeasibility certificates. The script also checks
the one-stage boundary proof's multiplier identities for $2\leq m\leq12$.
The JSON certificate contains exact fractions, not rounded solver output.

\texttt{scripts/verify\_model\_audit.py} retains the earlier exact audit,
including the literal model construction, the arXiv sign defect, and the
one-stage LP example. Its historical mode name \texttt{repaired} refers to
the arXiv signs; those signs are already correct in the journal.
\texttt{scripts/explore\_physical.py} implements a floating-point SCIP MILP
for $\BB$ and a nonconvex quadratic model for $\CC$. The small stage-gap and
initial one-stage experiments met their numerical expectations. The three
new counter-current potential-gadget runs reached their 20-second time
limits without settling the target values, including the YES instance;
their logs explicitly say FAIL for an unresolved numerical expectation.
They do not contradict or certify a theorem. Exact witnesses and the
algebraic proof replace those solver inferences.

The result settles the polynomial-versus-NP-hard question for the two
explicit physical completions, while the source's literal display is easy.
It does not assert hardness for every possible engineering superstructure,
for bounded numbers of streams, or for unrestricted stage counts chosen
without a supplied stage budget. General NP membership of $\CC$ remains
unproved here; strong NP-hardness does not depend on that membership.
The bounded novelty audit found no previous resolution, but its citation
coverage was incomplete and several citing works were screened only by
metadata or abstracts. Detailed sources and access limitations are recorded
in \texttt{references.md}. Independent verification of this preliminary
argument is pending.
\paperreferences
\end{document}
